% META: {"id":"1SPE-PRODSCAL-EX-045","chapitre":"1SPE-PRODUIT-SCALAIRE","type_objet":"exercice","capacites_codes":[],"extension_codes":["X1"],"programme_alignment":"OPTIONAL_EXTENSION","extension_label":"Approfondissement — Vers la Terminale","methodes":["M5"],"parcours":2,"competences":["calculer"],"duree_min":5,"mode_creation":"ex_nihilo","sources_inspiration":[],"fichier_tex":"chapitres/1SPE-PRODUIT-SCALAIRE/exercices/1SPE-PRODSCAL-EX-045.tex","corrige_tex":"chapitres/1SPE-PRODUIT-SCALAIRE/corriges/1SPE-PRODSCAL-CO-045.tex","status":"generated"}
% BEGIN-VERIFY
% from sympy import *
% # Median formula: AM^2 = (2*AB^2 + 2*AC^2 - BC^2)/4
% AM_sq = Rational(2*36 + 2*64 - 100, 4)
% assert AM_sq == 25
% assert sqrt(25) == 5
% END-VERIFY
\begin{exercice}{1SPE-PRODSCAL-EX-045}{2}{8}
\begin{center}\textbf{Approfondissement — Vers la Terminale}\end{center}
Dans un triangle $ABC$, $AB = 6$, $AC = 8$ et $BC = 10$.
\begin{enumerate}
  \item Calculer la longueur de la médiane issue de $A$ a l'aide de la formule de la médiane.
  \item Quelle est la nature du triangle ?
\end{enumerate}
\end{exercice}
